% chktex-file 44
\documentclass{article}
\usepackage[style=authoryear, backend=biber]{biblatex}
\usepackage[margin=2.54cm]{geometry}
\usepackage{parskip}
\usepackage{bbding}
\hyphenpenalty=1000000
\setlength{\parskip}{\baselineskip}
\addbibresource{ref.bib}
\setlength{\parindent}{0pt}
\renewcommand{\contentsname}{Table Of Contents}
\title{EGB320 Final Report: Mobility Subsystem}
\author{Darcy Stewart
\\N11977965
\\ Group 3}
\begin{document}

\maketitle
\pagebreak
\newgeometry{left=5cm, right = 5cm}
\begin{center}
   \tableofcontents 
\end{center}
\pagebreak
\restoregeometry{}
\section{Problem Description}
The problem set forward by the client was that shelves needed restocking, and it was believed that a robotic system would provide a more cost-effective method for shelf
restocking. To evaluate the suitability of a robotic warehousing system a Technology Readiness Level 3(TRL3) prototype. This prototype was designed to work in a scaled
down mock warehouse with shelves, items and obstacles that would simulate the major factors around navigating and moving stock around a grocery store.
\par
For this prototype some adjustments were made to allow for more streamlined development. The first major adjustment was to have each different element in the store be
represented by an object of a different color and shape as shown below:
\par
\begin{center}
\begin{tabular}{ |c|c| } 
 \hline
 Element & Representation \\ 
 \hline
 Obstacle & Green Colouring \\ 
 Item & Orange Colouring\\ 
 Aisle and Picking Station Marker & Black Colouring\\
 Shelves & Blue Colouring\\
 Picking Station & Yellow Colouring\\
 \hline
\end{tabular}
\end{center}
These adjustments were made to allow for a more simplified vision system which in turn would allow for rapid development. Further prototypes will build on this model to
allow for more complex vision processing. Along with this the shelf heights were the same for all aisles. Allowing for simplified item placement and navigation code as
no matter the bay the heights will always remain the same. The final simplification is the distance between the picking stations, this distance allows for a lower level
accuracy being required to pick up an item.
\par
With these simplifications in mind 4 separate subsystems were identified that would be designed to integrate to a final product. The intention behind this was to further
simplify the system through having each system capable of operating separately with methods of communication to allow for integration. This allowed for the 4 subsystems
to be designed in parallel and setbacks in one element of the robot would not necessarily set back the overall progress of the robot.
\par
Following this the requirements for each subsystem were identified, for mobility these requirements were as follows:
\begin{itemize}
    \item Turn to a bearing on a single point in a given direction
    \item Turn at a set speed in a given direction
    \item Move forward and backward at a set speed
    \item Move forward and backward a set distance
    \item Have a chassis that is capable of carrying all other subsystems
    \item Have a power management system that can power all other subsystems
\end{itemize}
With these requirements identified background research was conducted to determine the most suitable solutions.

\pagebreak
\section{Background Research}
Research was conducted into various methods for robotic mobility used in similar applications in industries. Four major methods of mobility were identified:
\begin{itemize}
    \item Wheels
    \item Tracks
    \item Legs
    \item Mecanum Wheels
\end{itemize}
An investigation into the pros and cons of each system was conducted to allow for a clear evaluation of each option's suitability. The results of this investigation were
compiled into the table seen below.

\begin{center}
    \begin{tabular}{|c|c|c|c|c|}
        \hline
        Feature & Wheels & Tracks & Legs & Mecanum Wheels \\
        \hline
        Can spin in place & \XSolidBold{} & \Checkmark{} & \Checkmark{} & \Checkmark{}\\
        Reliably stable without extra sensors & \Checkmark{} &\Checkmark{} &\XSolidBold{} & \Checkmark{}\\
        Requires 2 motors or less & \Checkmark{} & \Checkmark{} &\XSolidBold{} & \XSolidBold{}\\
        Can easily have speed measured in all directions & \Checkmark{} & \Checkmark{} & \XSolidBold{} & \XSolidBold{}\\
        Easily programmable &\Checkmark{} & \Checkmark{} & \XSolidBold{} & \XSolidBold{}\\
        \hline
    \end{tabular}
\end{center}
As shown in the table above it was decided that tracks would be investigated as the primary option for mobility. This was due to its ability to navigate varying terrain
with small obstacles, while maintaining a simple drive system to ensure a reliable outcome. The Treads also have a high amount of surface friction due to the large surface
area in contact with the ground. However, there were concerns about the reliability of tracks as consulting an expert on their usage in the field found that they can be
prone to slipping off the retainer system over time. This concern was minimal as the environment that the tracks were being deployed in varied greatly from the environment
of a warehouse.
\par
Once the research into the mobility system was done it was crucial to determine a method in which the motors will be controlled. Initially it was thought that the Raspberry Pi utilized for the vision and navigation subsystems would also be used to control the motors. This plan was found to have issues one research found that the
Raspberry Pi GPIO pins were only capable of 50mA total for the entire pin header. This was a major concern as the motors can easily draw over 5 times this amount under
a standard load.
\par
Further research was conducted and found that the majority of robots that involved a mobility subsystem made use of one or multiple motor controllers. These motor
controllers take input signals from a central control system and take power directly from the power management system to drive the motors. This bypasses the 50mA current
limit of the GPIO pins and allows for stronger motors to be used.
\par
Once this had been decided the next significant piece of research was to determine a power management system that would be capable of supplying the robot and all of its
peripherals with power. As the robot's main power supply was an 8.4V Lithium-ion battery it's voltage would need to be regulated down to the required 5V input for the
Raspberry Pi. Multiple options were considered for this regulator component, the first being a simple 5V linear regulator due to its simplicity for implementation.
This idea was moved past due to concerns that the with the efficiency of the regulator. Following this research was conducted into the feasibility of a switching regulator.
Research found that switching regulators are significantly more efficient in their power regulation, however they add high levels of noise to the circuit which can
affect signals in the system. This however wasn't deemed to be a major issue as there were no components that are highly susceptible to noise. With research conducted
into all elements required for the mobility system it was then time to begin the actual design and component selection for the robot.
\pagebreak
\section{Design}

\end{document}